\section{Intelligences}

A network that solves problems is an intelligence. To make that precise takes three short definitions, building on one another.

\begin{definition}[Problem]
A problem is the gap between the current state and a desired state, written $P = D - C$, the distance between where a system is and where it wants to be.
\end{definition}

\begin{definition}[Solution]
A solution is a path found from the current state to the desired one, a route through the space of configurations that arrives at the goal.
\end{definition}

\begin{definition}[Intelligence]
Intelligence is the capacity to find such paths, the standing ability to route from where a system is to where it wants to be, across many problems rather than one.
\end{definition}

Picture the space of everything a system could be as a vast network, each configuration a node, each move the law allows an edge between nodes. The current state is one node and the goal is another. A problem is the distance between them, a solution is a path that joins them, and intelligence is the knack for finding such a path quickly through a network far too large to search blindly. The mesh, with its solved addresses and its branching tree, is exactly such a network, which is why navigation and intelligence are the same skill seen twice.

None of this is only a definition. The problem-solving was measured. A search that knows its goal and follows the arrow toward it reaches a target a hundred steps away in about a hundred steps, where a blind walk would need the exponential of that, a gap of a billion to one by thirty steps, so intention is the difference between finding a thing and never finding it. Coordinated across a network, the bulk's logarithmic diameter lets a problem spread over many selves be solved in a number of hops that grows only as the logarithm of its size, where a flat arrangement would need the cube root. And a population of such solvers improves over generations, tuning how far ahead it plans to the hardness of its task. Intelligence is not only definable on the substrate, it is fast there.

A self has desires built in, the simplest being to persist, to stay bound against the bath. The gap between that desire and its present means is its first problem, and any structure that narrows the gap is acting intelligently, whether by learning a better model, by reaching a new capability, or by tempering a desire it cannot meet. Take a self whose body is being worn down faster than it can repair. It can intelligize the problem three ways, by pumping harder to raise repair, by shedding the vulnerable part to lower the demand, or by moving to a gentler region. Each closes $P$, and a system that can find such moves is intelligent in the plainest sense.

Networks are good at this, because they hold more model and try more strategies than any member alone. Intelligence is therefore a natural product of the tower, rising wherever the means to close a gap outrun any single self. This is intelligence in the working, engineering sense, problem-solving by a bounded system. It is not yet the cosmic intelligence the part on reality reaches for, the largest pattern in which the whole tower is one self. That is a separate and far bolder idea, and it is marked as such when it arrives. Here the claim is modest and mechanical. Wherever selves network, intelligence is what their networking is for.
